\subsection{整环的除性理想}

定义 1. 设 A 为整环，K 为其分式域。K 的每个满足下述条件的 A-子模 a 都称为 A 的分式理想（滥用语言时也称为 K 的分式理想）：存在 A 中元素 \(d \neq 0\)，使得 da ⊂ A。

 的每个有限生成 A-子模  都是分式理想：若 \(\mathfrak{a}\) 是 \(\mathbf{K}\)\((a_{i})_{1\leqslant i\leqslant n}\) 的一组生成元，可写成 \(\mathfrak{a}\)，其中 \(a_{i} = b_{i} / d_{i}\)、\(b_{i}\in \mathbf{A}\) 且 \(d_{i}\in \mathbf{A}\)；若 \(d_{i}\neq 0\)，显然 \(d = d_{1}\ldots d_{n}\)。特别地，\(da\subset A\) 的单生成 A-子模都是分式理想（回忆它们在代数第六章 § 1，第 5 节中被称为分式主理想）。若 \(\mathbf{K}\) 是 Noether 环，每个分式理想都是有限生成的 A-模。\(\mathbf{A}\) 的分式理想的每个 A-子模都是分式理想。\(\mathbf{A}\) 的每个理想都是分式理想；为避免混淆，也将它们称为 \(\mathbf{A}\) 的整理想。\(\mathbf{A}\) 的整理想。

记 I(A) 为 A 的非零分式理想所成的集合。给定 I(A) 的两个元素 a、b，将关系“每个包含 a 的分式主理想也包含 b”记作 （或 \(a \prec b\)）；显然此关系是 I(A) 上的预序。令 R 表示相伴的等价关系“\(b \succ a\) 且 ”（集合论，第三章 § 1，第 2 节），令 D(A) = I(A)/R 为商集；称 D(A) 的元素为 A 的除子，并且对每个分式理想 a  I(A)，将 a 在 D(A) 中的典范像记作 div a（或 div \(\prec b\)\(b \prec a\) a），称 div a 为 a 的除子；若 a = Ax 是分式主理想，则以 \(a \in I(A)\)\(_{A}\) 代替 \(\operatorname{div}(x)\)，并称 \(\operatorname{div}(\mathbf{A}x)\) 为 x 的除子；形如 \(\operatorname{div}(x)\) 的 D(A) 的元素称为主除子。取商后，I(A) 上的预序 \(x\)\(\operatorname{div}(x)\) 在 D(A) 上定义了一个序，记作 \(\prec\)。\(\leqslant\)。

对于每个 ，根据假设存在某个 \(\mathfrak{a} \in \mathrm{I}(\mathbf{A})\) 中的 ，使得 \(d \neq 0\)；因此，包含 \(\mathbf{A}\)\(\mathfrak{a} \subset \mathrm{Ad}^{-1}\) 的分式主理想的交  是 \(\tilde{\mathfrak{a}}\) 的元素。显然，关系 \(\mathfrak{a}\)\(\mathrm{I}(\mathbf{A})\) 等价于关系 \(\mathfrak{a} \prec \mathfrak{b}\)；因此关系 \(\tilde{\mathfrak{a}} \supset \tilde{\mathfrak{b}}\) 蕴含 \(\mathfrak{a} \supset \mathfrak{b}\)。\(\mathfrak{a} \prec \mathfrak{b}\) 的两个元素  模 \(\mathfrak{a}, \mathfrak{b}\) 等价的充要条件是 \(\mathrm{I}(\mathbf{A})\)\(\mathbb{R}\)。\(\tilde{\mathfrak{a}} = \tilde{\mathfrak{b}}\)。

定义 2. I(A) 中每个满足  的元素  称为 A 的除性分式理想。\(\mathfrak{a}\)\(\mathfrak{a} = \tilde{\mathfrak{a}}\) 称为 A 的除性分式理想。

换言之，除性理想恰是非零的非空分式主理想族的交。除性理想的每个非零交仍是除性理想。若 a 是除性理想，则对所有 \(x \in K^{*}\)，ax 也是除性理想，因为映射 \(b \mapsto bx\) 将分式主理想集双射到自身。对所有 \(a \in I(A)\)，\(\tilde{a}\) 是包含 a 的最小除性理想，且与 a 模 R 等价。此外，若 b 是与 a 模 R 等价的除性理想，则 \(\tilde{a} = \tilde{b} = b\)。因此 \(\tilde{a}\) 是满足 div a = div b 的唯一除性理想 b（换言之，将映射 \(a \mapsto div a\) 限制到除性理想集上所得映射是单射）。

设  与 \(\mathfrak{a}\) 为 \(\mathfrak{b}\) 的两个分式理想。回忆（第一章 § 2，第 10 节），\(\mathbf{K}\) : \(\mathfrak{b}\) 表示满足 \(\mathfrak{a}\) 的  所成的集合；这显然是一个 A-模；若 \(x \in \mathbf{K}\) 且 \(x\mathfrak{a} \subset \mathfrak{b}\)\(\mathfrak{b} \in \mathrm{I}(\mathrm{A})\)，则 \(\mathfrak{a} \in \mathrm{I}(\mathrm{A})\)；因为若 \(\mathfrak{b} \colon \mathfrak{a} \in \mathrm{I}(\mathrm{A})\) 是 \(d\) 中满足 \(\mathbf{A}\) 且 \(db \subset \mathbf{A}\) 的非零元素，且 \(da \subset \mathbf{A}\) 是 \(a\) 中的非零元素，则 \(\mathbf{A} \cap \mathfrak{a}\)；另一方面，若 \(da(\mathfrak{b} : \mathfrak{a}) \subset \mathbf{A}\) 属于 \(b \neq 0\)，则 \(\mathfrak{b}\)，故 \(bda \subset \mathfrak{b}\) : \(bd \in \mathfrak{b}\)，且 \(\mathfrak{a}\)。\(\mathfrak{b} : \mathfrak{a} \neq 0\)。

b: a 的定义还可写成：

\[
\mathfrak {b}: \mathfrak {a} = \bigcap_ {x \in \mathfrak {a}, x \neq 0} \mathfrak {b} x ^ {- 1}.\tag{1}
\]

命题 1. (a) 若 \(\mathfrak{b}\) 是除性理想且 \(\mathfrak{a} \in \mathbf{I}(\mathbf{A})\)，则 \(\mathfrak{b}\) : \(\mathfrak{a}\) 是除性理想。

\begin{enumerate}
\def\labelenumi{(\alph{enumi})}
\setcounter{enumi}{1}
\item
令 \(\mathfrak{a},\mathfrak{b}\) 属于 I(A)。要有 \(\operatorname {div}\mathfrak{a} = \operatorname {div}\mathfrak{b}\)，充要条件是 A: \(\mathfrak{a} = \mathbf{A}\) : \(\mathfrak{b}\)。
\item
对所有 \(\mathfrak{a} \in \mathbf{I}(\mathbf{A})\)，都有 \(\tilde{\mathfrak{a}} = \mathbf{A} : (\mathbf{A} : \mathfrak{a})\)。
\end{enumerate}

断言 (a) 由公式 (1) 立即得到，因为若  是除性理想，则对所有 \(\mathfrak{b}\)， 也是除性理想。\(\mathfrak{bx}^{-1}\)\(x \neq 0\) 也是除性理想。

为证明 (b)，令  表示包含 \(\mathrm{P}(\mathfrak{a})\) 的分式主理想所成的集合；关系 \(\mathfrak{a}\) 等价于 \(Ax \in \mathrm{P}(\mathfrak{a})\)，从而等价于 \(x^{-1}\mathfrak{a} \subset A\) : \(x^{-1} \in A\)。由于关系 \(\mathfrak{a}\) 按定义等价于 \(\operatorname{div} \mathfrak{a} = \operatorname{div} \mathfrak{b}\)，它也等价于 A : \(\mathrm{P}(\mathfrak{a}) = \mathrm{P}(\mathfrak{b})\) : \(A\)\(\mathfrak{a} = A\)。\(\mathfrak{b}\)。

最后，由于 ，有 。在此式中以  代替 ，可见 \(\mathfrak{a}(\mathbf{A}:\mathfrak{a})\subset \mathbf{A},\mathfrak{a}\subset \mathbf{A}:(\mathbf{A}:\mathfrak{a})\)\(\mathfrak{a}\)；另一方面，关系 \(\mathbf{A}:\mathfrak{a}\)\(\mathbf{A}:\mathfrak{a}\subset \mathbf{A}:(\mathbf{A}:(\mathbf{A}:\mathfrak{a}))\) 蕴含\(\mathfrak{a}\subset \mathbf{A}:(\mathbf{A}:\mathfrak{a})\) 蕴含

\[
\mathrm{A}: \mathfrak {a} \supset \mathrm{A}: (\mathrm{A}: (\mathrm{A}: \mathfrak {a})).
\]

因此 A: a = A : (A : (A : a))，由 (b) 得 div a = div(A : (A : a))。由于 A : (A : a) 根据 (a) 是除性理想，当然 \(\tilde{a} = A\) : (A : a)，这就证明了 (c)。

注。在上述证明过程中已经证明：对每个理想 \(a \in I(A)\)，都有 A: a = A: (A: (A: a))；这是集合论第三章 § 1，第 5 节命题 2 的特例。

命题 2. (i) 在 D(A) 中，每个非空有上界的集合都有最小上界。更确切地，若 \((\mathfrak{a}_t)\) 是 I(A) 中有上界的非空元素族，则

\[
\sup _ {i} (\operatorname{div} a _ {i}) = \operatorname{div} \left(\bigcap_ {i} \tilde {a} _ {i}\right).
\]

\begin{enumerate}
\def\labelenumi{(\roman{enumi})}
\setcounter{enumi}{1}
\tightlist
\item
在 D(A) 中，每个非空有下界的集合都有最大下界。更确切地，若 \((\mathfrak{a}_{\mathfrak{t}})\) 是 I(A) 中有下界的非空元素族，则
\end{enumerate}

\[
\inf _ {i} (\operatorname{div} a _ {i}) = \operatorname{div} \left(\sum_ {i} a _ {i}\right).
\]

\begin{enumerate}
\def\labelenumi{(\roman{enumi})}
\setcounter{enumi}{2}
\tightlist
\item
集合 D(A) 是一个格。
\end{enumerate}

设 \((\mathfrak{a}_i)\) 是 I(A) 中有上界的非空元素族。说除性理想 \(\mathfrak{b}\) 是这个族的上界，就是说它包含于所有 \(\tilde{\mathfrak{a}}_i\) 中，也就是说 \(\mathfrak{b}\) 包含于 \(\bigcap_{i} \tilde{\mathfrak{a}}_i\) 中。因此 \(\bigcap_{i} \tilde{\mathfrak{a}}_i \neq (0)\)，且 \(\bigcap_{i} \tilde{\mathfrak{a}}_i\) 是除性理想，这就证明了 (i)。

现在令 \((\tilde{a}_t)\) 为 I(A) 中有下界的非空元素族。说除性理想 \(b\) 是这个族的下界，就是说它包含所有 \(\tilde{a}_t\)，也就是说（因为 \(b\) 是除性理想）它包含所有 \(a_t\)，或者说 \(b \supset \sum_{i} a_i\)。这就证明了 (ii)。

最后，为证明 (iii)，根据 (i) 与 (ii)，只需证明当 a、b 属于 I(A) 时，集合 \(\{a, b\}\) 在 I(A) 中既有上界又有下界；现在它以上界 \(a \cap b\) 为界（该集合不同于 (0)）。它以下界 \(a + b\) 为界，因为 \(a + b \in I(A)\)：若 d 与 \(d'\) 是 A 中满足 \(da \subset A\) 且 \(d'b \subset A\) 的非零元素，则 \(dd'(\mathfrak{a} + \mathfrak{b}) \subset \mathbf{A}\)。

推论。若 \(x, y\) 与 \(x + y\) 属于 \(\mathbf{K}^*\)，则 \(\operatorname{div}(x + y) \geqslant \inf(\operatorname{div}(x), \operatorname{div}(y))\)。\(\mathbf{A}(x + y) \subset \mathbf{A}x + \mathbf{A}y\)，从而 \(\operatorname{div}(x + y) \geqslant \operatorname{div}(\mathbf{A}x + \mathbf{A}y)\)。
% semantic-fragments: legacy-input-seam
\relax{}
